\documentclass[conference]{IEEEtran}
\usepackage[utf8]{inputenc}
\usepackage[T1]{fontenc}
\usepackage{cite}
\usepackage{amsmath,amssymb}
\usepackage{graphicx}
\usepackage{booktabs}
\usepackage{url}

\title{Measuring Validator Utility in Generate--Verify--Repair NetOps Pipelines: a Confirmatory Paired Study with Shared Initial Candidates}

\author{
\IEEEauthorblockN{Autonomous AI Researcher}
\IEEEauthorblockA{CNSM 2026 Agentic AI Researcher Track}
}

\begin{document}

\maketitle

\begin{abstract}
We preregistered and executed a paired confirmatory study (AC-2026-03) to evaluate whether a multidimensional validator metric suite predicts end-to-end agentic NetOps success better than a single verifier pass rate. Using a synthetic intent\ensuremath{\rightarrow}configuration generator and a hosted generate--validate--repair (GVR) adapter under shared\_initial\_candidate semantics, we ran 200 paired tasks (one shared initial generation per task reused for baseline and guarded). The deterministic validator (deterministic\_netops\_validator\_v5, v5.0) judged every sampled initial candidate valid; no repairs were applied. Baseline = 200/200, guarded = 200/200, paired difference = 0.00 (bootstrap 95\% CI [0.00, 0.00], exact McNemar two-sided p = 1.0; bootstrap\_seed = 1729, resamples = 10000). The observed ceiling invalidated the preregistered power assumptions (targeted 10 percentage-point paired improvement) and rendered the primary confirmatory test uninformative for the originally planned effect. We report preregistered design choices, precise execution accounting, representative validator-trace excerpts, quantitative analysis, failure-mode diagnostics, and artifact-grounded prescriptions for follow-up preregistered experiments (fault injection, multi-sample generation, multi-validator comparison) required to estimate validator coverage and repair value.
\end{abstract}

\section{Introduction}
Automating NetOps workflows with generative models requires reliable mechanisms to avoid harmful, silent, or wrong-but-plausible actions. A common operational pattern is generate--verify--repair (GVR): a generator (LLM) proposes a candidate configuration or plan, a deterministic validator mechanically checks correctness and policy constraints, and, when necessary, a repair module synthesizes corrected candidates or triggers human review. Prior demonstrations of GVR-like architectures in code and systems motivate their adoption in NetOps, but systematic, preregistered measurements of validator coverage (false-positive/false-negative rates), detection latency, and predictive usefulness for end-to-end guarded success are scarce [1].

We preregistered study AC-2026-03 with a paired design that isolates the effect of downstream validator-triggered repair under shared\_initial\_candidate semantics: for each task the generator produced a single initial candidate that was reused by both the baseline (no repair) and guarded (validator+repair permitted) conditions. Under these semantics any paired difference can only arise downstream from validator-triggered repair, not from independent generator sampling. The primary estimand was the paired\_success\_rate\_difference\_guarded\_minus\_baseline and the preregistered confirmatory hypothesis (H1) expected that a multidimensional validator-metric suite would explain more variance in end-to-end success than a single verifier pass rate. Our main contribution is a fully preregistered, artifact-archived experiment and an exact, transparent report of a degenerate confirmatory outcome (ceiling) with concrete, artifact-grounded prescriptions for follow-up designs that can measure validator utility.

\section{Related Work}
Work across code, systems, and NetOps communities has proposed and prototyped generate--verify--repair or closed-loop verifier architectures to reduce stochastic generator error and enable deterministic acceptance criteria [1]. Recent NetOps and AIOps benchmark efforts document substantive operational failure modes (ticket noise, false premises, and wrong-but-plausible outputs) that complicate end-to-end automation and motivate insertion of verifier layers [2],[3]. Independent system reports and preprints propose self-healing orchestrators, auditable decision traces, and verifier-guided repair to reduce silent or action-triggering failures in tool-augmented LLM systems [4]. Surveys of LLMs for telecommunications and NetOps discuss adaptation and verification trade-offs and motivate metricized evaluation of verification cost versus nominal correctness [5].

These verified records motivate a measurement focus on validator soundness and operational impact but do not provide a standardized, empirically estimable validator-metric suite tailored to NetOps GVR flows. The present preregistered experiment was designed to begin filling that measurement gap under strict preregistration and archived-execution practices; when interpreted against the cited literature we treat prior demonstrations as architectural and motivational rather than as establishing empirical coverage for NetOps validators.

\section{Methodology}
Preregistered design and confirmatory intent. We preregistered study AC-2026-03 as a paired confirmatory experiment to isolate downstream validator/repair effects from generator variability. The registry specified adapter\_family = hosted\_netops\_gvr\_v1 and generation\_semantics = shared\_initial\_candidate with initial\_generation\_calls\_per\_task = 1. The primary estimand was the paired\_success\_rate\_difference\_guarded\_minus\_baseline; the primary test was an exact two-sided McNemar with paired bootstrap percentile confidence intervals (bootstrap\_resamples = 10000, bootstrap\_seed = 1729). The design targeted detection of an absolute 0.10 paired improvement (\ensuremath{\alpha} = 0.05, \textasciitilde{}80\% power) under conservative planning assumptions and a preregistered sample size of task\_count = 200 pairs (planned\_episode\_count = 400).

Task generation and constraints. Tasks were drawn deterministically from deterministic\_netops\_task\_generator\_v5 (task\_family = intent\_configuration\_repair\_v1). Each task manifest entry contains: an initial\_state, expected\_final\_state, a human-readable intent string, allowed\_fields and allowed\_touched\_fields, explicit workflow\_policies (e.g., minimum\_active\_in\_groups, must\_be\_down\_for\_fields), and a required\_command\_sequence. Task selection followed the preregistered deterministic index rule challenging\_workflow\_stress\_v1 to produce a stress-profiled suite (levels labeled 'very\_hard' and 'extreme' in difficulty fields). Contamination and exclusion rules were preregistered: identical SHA-256 duplicates are excluded and replaced deterministically; near-duplicate tasks (embedding cosine > 0.95) are deterministically excluded and replaced. The execution/analysis artifacts (analysis/contamination\_summary.csv) record no contamination exclusions.

Model, sampling, and provider accounting. The declared model in execution/model\_configuration.json is "gpt-5-mini" and the provider bundle is recorded as "openai\_responses". Preregistered model\_scope = gpt-5-mini and maximum\_model\_calls = 400. The adapter executed exactly one shared initial generation per task (model\_calls\_used = 200, as recorded in the execution manifest). Important artifact-grounded sampling detail: execution/model\_configuration.json records temperature = null (temperature = null/unset). Explicitly, null/unset is not evidence of deterministic hosted-model sampling and does not imply temperature = 0; provider-side sampling behavior may remain stochastic and is captured in provider-call audit fields. Deterministic\_reconciliation documents hosted\_model\_call\_event\_count = 200 and stage\_counts.shared\_initial\_generation = 200, confirming the single shared generation per pair semantics.

Validator, scoring rule, and repair policy. The pipeline used deterministic\_netops\_validator\_v5 (validator\_version = 5.0) as the mechanistic validator. The validator implements rule-based intent-constraint checks and emits per-episode fields used by scoring: normalized\_configuration, parsed\_command\_count, required\_command\_count, observed\_touched\_field\_count, violation\_count and violation\_codes, valid (boolean), repair\_applied (boolean), and repair\_provider\_trace fields when a repair is invoked. The preregistered binary scoring rule full\_intent\_constraint\_validation maps validator.valid (true/false) to the confirmatory success label. The pipeline permitted at most one repair call when the validator rejected a candidate (repair budget per episode = 1); repair events would be logged with repair\_applied = true and repair\_provider\_trace populated. In practice, the archive shows repair\_applied = false for all episodes and null repair trace fields.

Execution-level controls, exclusions, and missingness. Operational execution constraints were preregistered: maximum\_attempts\_per\_call = 1 (no manual retries), deterministic replacement indices for excluded tasks, and a complete-pair primary analysis rule (paired\_binary\_analysis\_v1 uses only complete pairs). The manifested execution ran to completion with completed\_episode\_count = 400, failed\_episode\_count = 0, and complete\_pair\_count = 200. Missingness artifacts (analysis/missingness\_summary.csv) show no baseline-only, guarded-only, or both-missing pairs. All provider-call, validator-trace, and raw-response artifacts were archived for reproducibility.

Analysis plan and inferential checks. The primary confirmatory test was the exact McNemar two-sided procedure applied to the paired contingency (n\_11, n\_10, n\_01, n\_00) with bootstrap percentile 95\% confidence intervals (10,000 resamples, seed = 1729) for the paired success-rate difference. Secondary preregistered analyses included: per-validator FP-rate and FN-rate estimation (averaged across tasks), mean detection latency per validator, and a linear model predicting a simulated operational-impact score from the validator metric vectors (report adjusted R\textasciicircum{}2). Multiplicity and exploratory analyses were declared: the primary estimand was unique and exploratory subgroup/error-class analyses would control false discovery rate when multiple p-values were reported. The preregistration also specified sensitivity checks: treating missing or incomplete episodes as failures (zero) and bootstrap stability checks (1,000 resamples recommended during planning; final confirmatory bootstrap used 10,000 resamples as recorded).

Why these choices. The paired shared\_initial\_candidate design was chosen to provide a clean scientific isolation: any improvement observed in the guarded condition must arise from deterministic validator detection and repair (or from repair-induced acceptance differences), not from independent generator sampling differences. The preregistered contamination rules, single-attempt policy, and deterministic task indexing were selected to reduce confounds and keep the confirmatory estimand interpretable. The trade-off is explicit and documented: shared\_initial\_candidate deliberately suppresses generator variability and therefore increases the risk of ceiling/floor outcomes when the single sampled candidate systematically agrees with the validator, a risk that materialized in the present run.

Archive and artifact linkage for reproducibility of methods. All method-level configuration objects and runtime traces referenced above are present in the archived execution bundle: execution/model\_configuration.json (declared\_model\_version = "gpt-5-mini", temperature = null, maximum\_attempts\_per\_call = 1), execution/task\_manifest.jsonl (deterministic\_netops\_task\_generator\_v5 payloads and required sequences), execution/raw\_results.jsonl (raw model responses), execution/scoring/*.json (per-episode validator traces), and analysis artifacts (analysis/paired\_contingency\_table.csv, analysis/condition\_summary.csv, analysis/missingness\_summary.csv). The preregistration SHA-256 is recorded in the execution manifest (preregistration\_sha256 = de37b485\allowbreak{}421cb089\allowbreak{}5a98df38\allowbreak{}b6b3baad\allowbreak{}1dc7c13c\allowbreak{}81065e3d\allowbreak{}8240d2a9\allowbreak{}4a0ff844). These artifacts support independent verification of method-level choices, confirm that the single-sample-per-task and shared\_initial\_candidate semantics were enforced, and document the validator interface used to produce the binary success labels.

\section{Results}
Execution accounting and raw outcome summary. The run completed exactly per preregistration: planned\_episode\_count = 400, completed\_episode\_count = 400, failed\_episode\_count = 0, complete\_pair\_count = 200, and model\_calls\_used = 200 (one shared initial generation per pair). Provider-call audit entries show hosted\_model\_call\_event\_count = 200, completed\_hosted\_model\_call\_event\_count = 200, cache\_hit\_count = 0, and cache\_miss\_count = 200. The analysis condition summary (analysis/condition\_summary.csv) reports: baseline successes = 200, baseline failures = 0 (success\_rate = 1.0); guarded successes = 200, guarded failures = 0 (success\_rate = 1.0).

Primary confirmatory statistics and exact contingency. The paired contingency (analysis/paired\_contingency\_table.csv) is n\_11 = 200, n\_10 = 0, n\_01 = 0, n\_00 = 0; that is, every pair was valid under both baseline and guarded. The paired\_success\_rate\_difference\_guarded\_minus\_baseline equals 0.00. The preregistered exact McNemar two-sided test returned p = 1.0 because there are zero discordant pairs (n\_10 + n\_01 = 0). The preregistered bootstrap-percentile confidence interval (bootstrap\_resamples = 10000, bootstrap\_seed = 1729) is [0.00, 0.00], reflecting a point mass at zero in the resampled paired-difference distribution. These analysis outputs are recorded verbatim in analysis artifacts (analysis/analysis\_log.jsonl, analysis/paired\_contingency\_table.csv).

Representative validator-trace and why no repairs occurred. A typical archived validator\_trace (pair-000004) exhibits the mechanistic checks the validator enforces and illustrates why no repair was triggered:

pair\_id: pair-000004
  normalized\_configuration: "interface edge2 admin down\textbackslash{}ninterface edge2 vlan 40\textbackslash{}ninterface edge2 admin up\textbackslash{}ninterface edge1 admin down"
  parsed\_command\_count: 4
  required\_command\_count: 4
  observed\_touched\_field\_count: 3
  valid: true
  repair\_applied: false
  validator\_id: deterministic\_netops\_validator\_v5
  validator\_version: 5.0

Equivalent validator\_trace.validation\_after.valid = true and repair\_applied = false hold for every archived episode; repair\_provider\_trace fields are null for all episodes. Representative task manifests show required\_command\_count equal to parsed\_command\_count for the accepted candidates, indicating structural completeness under the validator's rules. Because baseline and guarded reused the identical shared\_initial\_candidate per pair, these per-episode validator outputs explain the degenerate paired outcome: there were no validator-detected mechanical defects to correct.

Expanded diagnostics: what the artifacts reveal about failure modes and what remains unobserved. Several archived artifacts together explain the ceiling outcome without invoking unrecorded assumptions:
  - deterministic\_reconciliation and model\_configuration.json show the single-shared-generation accounting (one initial draw per task; model\_calls\_used = 200). This sampling regime eliminated generator-sampling variance as a potential source of discordant pairs.
  - analysis/contamination\_summary.csv and analysis/missingness\_summary.csv both report no contamination flags and complete\_pairs = 200, confirming the confirmatory dataset was contamination-free and complete as preregistered.
  - validator traces uniformly report valid = true and repair\_applied = false; thus per-validator FP/FN estimates are degenerate (no false positives and no false negatives observed relative to the validator's own criteria on this sampled set), and detection-latency/repair-invocation distributions are uninformative because no repair events occurred.

Power and inferential consequences in artifact terms. The preregistered power calculation targeted detection of an absolute paired improvement of 0.10 with \ensuremath{\alpha} = 0.05 and \ensuremath{\approx}80\% power under conservative assumptions (see preregistration). Those calculations assumed non-degenerate baseline variability (for example baseline p values in \{0.3,0.4,0.5,0.6,0.7\}). The artifact-grounded observed baseline\_success\_rate = 1.0 (200/200) makes the preregistered detectable effect logically unattainable in this execution: with no observed failures in baseline there are no discordant baseline-only episodes to be corrected by the guarded condition, so the paired estimator cannot register the preregistered 10-percentage-point improvement. We therefore report the confirmatory null outcome as an artifact-valid, design-driven ceiling rather than as evidence that validators lack operational value generally.

What the archived traces allow and what they do not. The execution bundle contains full per-episode normalized\_configurations, required\_command\_count and parsed\_command\_count fields, provider-call traces, and analysis CSVs. These artifacts enable reproducible reanalyses under alternative scoring rules (for example, scoring that treats specific semantic mismatches as failures) or under changed sampling semantics (multi-sample per task) or injected faults. What the current artifacts do not support is estimation of repair effectiveness or validator false-negative rates in unperturbed production-like inputs because repair\_applied = false for every episode and no injected faults were executed in this run.

Concluding diagnostic summary. The degenerate all-valid outcome arises from the conjunction of (1) a deterministic task generator that produced candidates aligned with the validator's mechanistic constraints for the sampled seeds, (2) a single-sample-per-task shared\_initial\_candidate design that removed generator variability as a potential source of discordance, and (3) a deterministic validator that enforces structural and policy constraints but does not attempt to model operator intent beyond the encoded workflow\_policies. These artifact-grounded determinants fully account for the lack of observed repairs and the uninformative confirmatory statistic; follow-up preregistered arms (fault injection, multi-sample generation, multi-validator comparison) are required --- and are supported by the archived infrastructure --- to produce non-degenerate estimates of validator coverage, FP/FN rates, repair invocation utility, and detection latency.

\section{Discussion}
Design implications of shared\_initial\_candidate semantics. The preregistered shared\_initial\_candidate choice isolates downstream validator effects by ensuring both conditions evaluate the same initial candidate for each pair. This isolation is scientifically valuable because any observed paired difference must derive from validator-triggered repair. It also concentrates the risk of a single-sample ceiling: if the generator sample aligns with the validator for all tasks, the design yields zero observed variance in the primary estimand, as occurred here. We emphasize this property because it determines what the paired estimator can and cannot measure: it cannot detect differences attributable to generator sampling or to alternative prompt constraints when the sampling semantics were shared.

Operational interpretation and validator scope. The deterministic validator (deterministic\_netops\_validator\_v5) reliably enforces mechanistic intent-constraint checks (parsed\_command\_count, required\_command\_count, violation\_codes, valid flag). However, artifact-grounded evidence here shows that correctness-by-structure does not imply semantic intent alignment under operator expectations. That limitation is intrinsic to validators that lack an external intent model or operator feedback; detecting semantic mismatches requires either richer intent models, human-in-the-loop signals, or tailored semantic validators. The present execution does not imply validators are unnecessary---rather, it demonstrates that in a synthetic bank and single-sample regime the validator had nothing mechanical to correct.

Prescriptions for follow-up preregistered experiments (artifact-grounded). The archived infrastructure and artifacts directly support concrete next-step preregistered arms that would produce estimable validator FP/FN and repair utility while preserving preregistration rigor:
  1) Fault-injection arm (preregistered): deterministically inject defined fault classes into a controlled fraction f of reference answers (syntactic errors, configuration-logic faults, semantic-intent swaps). This increases baseline invalidity and permits measurement of repair invocation rates and post-repair success.
  2) Multi-sample generation: set initial\_generation\_calls\_per\_task = k >= 3 independent draws per task with independent seeds recorded. Allow the guarded pipeline to accept any validator-approved draw or to perform repair. This reintroduces generator variability and creates opportunities for validator rejections and repairs.
  3) Multi-validator and multi-model comparison: add alternative validator implementations and at least one additional model family. This will quantify sensitivity to validator design and model diversity and improve external validity.

Each prescription is directly supported by the archived adapter semantics and task-generation determinism; implementing them preregistered would produce non-degenerate data suitable for estimating per-validator FP/FN, repair utility, detection latency, and predictive R\textasciicircum{}2 against an operational-impact proxy as originally planned.

Threats to validity. Internal validity: shared\_initial\_candidate intentionally isolates downstream effects but increased single-sample ceiling risk. Construct validity: validator.valid maps to mechanistic correctness but may fail to capture semantic intent misalignment. External validity: synthetic task bank, single-model (gpt-5-mini), and a single deterministic validator limit generalization to heterogeneous production networks and additional model families. Conclusion validity: the degenerate outcome yields zero variance in the primary estimand so inferential conclusions about validator benefit are not possible from this execution alone. These limitations are recorded in the execution and analysis artifacts and motivate the artifact-grounded follow-up arms described above.

\section{Limitations}
\begin{itemize}
\item Shared\_initial\_candidate semantics constrained inference: baseline and guarded reused the same initial candidate per pair; any paired difference could only arise from validator-triggered repair (not from independent generator sampling).
\item Synthetic task bank (difficulty 'very\_hard'/'extreme') is not equivalent to heterogeneous production networks (multi-vendor devices, real ticket noise); external validity is limited.
\item Single-model experiment: only gpt-5-mini was executed; no multi-model generality claims are made.
\item No repairs were applied because all initial candidates passed the deterministic validator; therefore the study could not estimate repair effectiveness or validator false-negative behavior in this execution.
\item Validator scope: deterministic\_netops\_validator\_v5 enforces mechanistic intent-constraint checks but does not guarantee detection of semantic intent misalignment (correct-by-structure but incorrect-by-intent).
\item Recorded execution/model\_configuration.json temperature = null/unset does not imply deterministic sampling; provider-side sampling behavior may vary and is documented in execution manifests.
\end{itemize}

\section{Conclusion}
We executed a fully preregistered paired confirmatory study (AC-2026-03) under shared\_initial\_candidate semantics to measure validator utility in NetOps GVR pipelines. For the synthetic task bank, the sampled gpt-5-mini generations, and deterministic\_netops\_validator\_v5 used here, every sampled initial candidate passed the validator's mechanistic checks so no repairs were applied and guarded success equaled baseline (200/200). The preregistered power assumptions (targeted 10-percentage-point improvement) were invalidated by this ceiling outcome. The result is artifact-grounded and reproducible from the archived bundle. We publish the exact execution and analysis artifacts to support replication and provide concrete, preregistered experiment prescriptions (fault injection, multi-sample generation, multi-validator comparison) that are required to produce estimable validator FP/FN behavior and repair value in operationally meaningful conditions.

\section*{Disclosure Statement}
Disclosure (concise): Declared model and provider: gpt-5-mini via provider bundle 'openai\_responses'. Orchestration/adapter: hosted\_netops\_gvr\_v1. Deterministic validator: deterministic\_netops\_validator\_v5 (validator\_version = 5.0). Preregistration: AC-2026-03 (preregistration\_sha256 = de37b485\allowbreak{}421cb089\allowbreak{}5a98df38\allowbreak{}b6b3baad\allowbreak{}1dc7c13c\allowbreak{}81065e3d\allowbreak{}8240d2a9\allowbreak{}4a0ff844). Master prompt immutable reference: provenance/master\_prompt.txt, SHA-256 = 1872df1e\allowbreak{}1805d2d9\allowbreak{}6940456c\allowbreak{}a016bd66\allowbreak{}5d1d5196\allowbreak{}add77f5a\allowbreak{}cdf1582b\allowbreak{}b39b15ba. Autonomy boundary: a human Research Director selected the broad topic family and approved the master prompt before the autonomy lock; after the lock the AI system autonomously performed literature review, preregistration, experiment design, execution, analysis, internal peer review and manuscript drafting without manual editing of technical content. Sampling/generation semantics: shared\_initial\_candidate (one initial sampled generation per task reused for baseline and guarded); execution/model\_configuration.json recorded temperature = null/unset (null/unset is not evidence of deterministic sampling and does not imply temperature = 0). Call budget and usage (preregistered): maximum\_model\_calls = 400; actual model\_calls\_used = 200 (one shared initial generation per pair). All raw outputs, per-episode validator traces, execution manifests, and analysis artifacts are archived in the supplied execution bundle (see execution/execution\_manifest.json). No human manually edited the manuscript technical content after the autonomy lock.

\begin{thebibliography}{99}
\bibitem{ref1} [1] R. Cadenas, "Convergent Correctness in Stochastic Code Generation: A Generate-Verify-Repair Architecture for Deterministic Validation of Autonomous AI Coding Agents in Regulated Environments," SSRN Electronic Journal, 2026. doi:10.2139/ssrn.6754899.
\bibitem{ref2} [2] K.-H. Tseng, N. Bogahawatta, Y. Ginige, K. Patel, K. Dakic, S. Seneviratne, "FaulT-Bench: Towards Benchmarking Network Troubleshooting LLM Agents under Unreliable User Tickets," arXiv:2608.27021, 2026.
\bibitem{ref3} [3] Y. Liu, C. Pei, L. Xu, B. Chen, M. Sun, Z. Zhang, Y. Sun, S. Zhang, K. Wang, H. Zhang, J. Li, G. Xie, X. Wen, X. Nie, M. Ma, P. Dan, "OpsEval: A Comprehensive IT Operations Benchmark Suite for Large Language Models," arXiv:2310.07637, 2023.
\bibitem{ref4} [4] R. S. Babu and A. Agrawal, "Self-Healing Agentic Orchestrators for Reliable Tool-Augmented Large Language Model Systems," arXiv:2606.01416, 2026.
\bibitem{ref5} [5] H. Zhou, C. Hu, Y. Yuan, Y. Cui, Y. Jin, C. Chen, H. Wu, D. Yuan, L. Jiang, D. Wu, X. Liu, J. Zhang, X. Wang, J. Liu, "Large Language Model (LLM) for Telecommunications: A Comprehensive Survey on Principles, Key Techniques, and Opportunities," IEEE Commun. Surveys \& Tutorials, 2024. doi:10.1109/COMST.2024.3465447.
\end{thebibliography}

\end{document}
